\begin{abstract}
Quartic and sextic centrifugal distortion constants are usually handled as
representation-dependent parameters of Watson's effective rotational
Hamiltonian, whether they are obtained from spectroscopic fits or from direct
quantum-chemical calculations.  In this work we reformulate the standard
quartic and standard sextic sectors in a common tensorial framework in which
Watson constants are treated as projected coordinates of underlying tensor
representatives reconstructed by Moore--Penrose pseudoinversion, which thus
acts as a natural gauge fixing of the inverse problem.

Within this formulation, standard quartic and standard sextic centrifugal
distortion are shown to belong to the same linear representation-equivalent
sector.  The key structural point is that the standard sextic geometry sector
can be written in the same first-level tensorial language that governs the
standard quartic problem, so that second derivatives of the inertia tensor need
not be treated as primitive objects.  The explicit anharmonic sextic terms then
enter only through the cubic force constants.  Representation and reduction
changes are therefore performed by lifting Watson constants to tensor
representatives, transporting the tensors under axis permutations, and
reprojecting them onto the target Hamiltonian parametrization.

This common standard sector provides the natural baseline from which the first
departure from linear representation equivalence can be identified.  On the
quartic side, the channel $H_{22}$ is shown to be the first contribution that
requires a genuinely new intrinsic tensorial object beyond the standard
first-level structure.  On the sextic side, the corresponding first
post-standard candidate is the term linear in $\tau(H_{22})$ inside the sextic
geometry functional.  These quantities are not introduced as complete
higher-order corrections, but as low-cost tensorial diagnostics of both the
sensitivity of fitted constants to representation/reduction choices and the need
to move beyond the standard perturbative sector.  The main text is therefore
focused on linear representation transport and its first breakdown, while the
detailed tensor constructions and derivations that justify this structure are
collected in the appendices.
\end{abstract}
